\chapter{Results and Discussions}
\label{sec:results_and_discussions}

\section{Illustrative Results}
All values in Table~\ref{tab:results} are invented for layout inspection. They
must not be interpreted as measured performance. In this numerical example,
the variant has a larger mean in both groups, but the difference is larger
for short requests. Because the groups have equal size in the example design,
the overall means are the simple averages of the two group means.

\begin{table}[htbp]
  \centering
  \caption{Invented scores used only to demonstrate results formatting}
  \label{tab:results}
  \begin{tabular}{lrrr}\hline
    Query group & Baseline & Variant & Difference\\\hline
    Short requests & 0.420 & 0.460 & +0.040\\
    Detailed requests & 0.380 & 0.390 & +0.010\\
    Overall & 0.400 & 0.425 & +0.025\\\hline
  \end{tabular}
\end{table}

A real discussion would report the uncertainty of the comparison and connect
it to the research questions. It should avoid describing a numerical increase
as a reliable improvement without the evidence needed for that statement.
The size of a change also needs practical interpretation: a difference may
be observable in an offline score without noticeably changing the usefulness
of the system for its intended users.

% Demonstration-only break: keep the image with discussion on a continuation page.
\clearpage
\section{Figure Example}
Figure~\ref{fig:logo} demonstrates a centred image with its caption below it.
The image is the existing template logo and conveys no experimental result.
In a completed thesis, replace it with a relevant figure and explain what the
reader should learn from it in the surrounding text.

\begin{figure}[H]
  \centering
  \includegraphics[width=5cm]{waseda_logo.pdf}
  \caption{Waseda University logo, included as a formatting example}
  \label{fig:logo}
\end{figure}

\section{Interpreting Differences}
An aggregate result can conceal different query-level patterns. Some queries
may benefit from a reranking stage while others lose a relevant document near
the top of the list. Inspecting the paired scores helps identify these cases,
but the interpretation should be grounded in the retrieved material rather
than inferred from the sign of the score difference alone.

The discussion should distinguish observations from possible explanations.
For example, a query might contain terminology that appears in several
unrelated contexts. A retrieved document can reveal such an ambiguity, but
that single case does not establish a general failure mode. Additional cases
or a targeted experiment would be needed to support a broader claim.

\section{Limitations}
The illustrative design has several limitations. A fixed collection may not
reflect changes in language or document availability. Incomplete relevance
judgements may leave some useful documents unrecognised by the evaluation.
A small query set may also give an unstable picture of group-level behaviour.
Each limitation affects the scope of the conclusion in a different way.

A useful limitations section explains those effects rather than merely listing
potential problems. It should identify which conclusions remain supported,
which are tentative and which would require a new experiment. Since this
chapter contains no real observations, its only conclusion is about how a
results discussion can be structured. The following chapter preserves that
boundary when summarising the example.
